\section{Introduction: The Missing Comparison}
\label{sec:intro}

A redistricting algorithm that produces compact, population-balanced maps
is useful only if we can say something about where it falls in the space
of all reasonable maps.  Without this comparison, a court or legislature
cannot distinguish between a map that is geographically sensible and one
that happens to be compact while achieving an extreme partisan outcome.

The ensemble literature provides exactly this comparison for human-drawn and
court-drawn maps.  \citet{herschlag2020quantifying} showed that the enacted
North Carolina congressional plan fell at the 99th percentile of Republican
seat share in a \recom{} ensemble of 24,000 plans---a result central to the
\case{Common Cause v.\ Rucho} litigation.  \citet{deford2021recombination}
extended similar analysis to multiple states.  But neither paper compared
algorithmic plans to their ensembles.

This paper fills that gap.  We place the \ar{} plan
\citep{dellaLibera2026b11} within published ensemble distributions for six
states.  These six states were chosen because:

\begin{enumerate}
  \item Published ensemble results are available from \citet{herschlag2020quantifying}
    or \citet{deford2021recombination};
  \item The states span a range of congressional delegation sizes ($k = 8$ to
    $k = 52$);
  \item The states include both competitive states (Wisconsin, Pennsylvania)
    and geographically sorted states (Georgia, Texas).
\end{enumerate}

\paragraph{Main finding.}
The \ar{} plan falls in the $61$st--$72$nd percentile of compactness and
near the $50$th percentile of partisan outcomes for five of six studied states.
Georgia is the outlier, with the \ar{} plan at the $38$th percentile of
Democratic seats, reflecting the state's strong geographic partisan sorting
rather than any algorithmic bias.

\paragraph{Legal implication.}
A plan at the $68$th PP percentile with \emph{deterministic, verifiable
provenance} is legally stronger than any ensemble plan, even one closer to
the partisan median, because it eliminates one entire class of challenge:
the challenge that the map was selected to achieve a partisan end.  When
the map is determined by the prime factorization of the seat count and the
METIS partition of the census graph, this challenge cannot succeed.

\paragraph{Organisation.}
Section~\ref{sec:states} describes the six states and data sources.
Section~\ref{sec:method} details the methodology.
Sections~\ref{sec:nc}--\ref{sec:tx-ca} present state-by-state results.
Section~\ref{sec:summary} provides the consolidated summary table.
Section~\ref{sec:legal} develops the legal argument.
Section~\ref{sec:conclusion} concludes.
